(1) 正入射往返相位 \(\delta=4\pi nL/\lambda=4\pi nL\nu/c\)。相继透射振幅构成公比 \(Re^{i\delta}\) 的级数，因两镜无吸收且对称，归一化强度 \(T(\nu)=(1-R)^2/|1-Re^{i\delta}|^2\)。化简得 Airy 公式 \(T=[1+F\sin^2(\delta/2)]^{-1}\)，其中 \(F=4R/(1-R)^2\)，峰值为一，这里忽略了镜相位的常数偏移。
(2) 谐振需 \(\delta=2\pi m\)，故 \(2nL=m\lambda\) 或 \(\nu_m=mc/(2nL)\)；若 \(n\) 不随频率改变，邻峰间距 \(\Delta\nu_{\rm FSR}=c/(2nL)\)。峰数 \(m\) 必须是整数，\(\Delta\nu_{\rm FSR}\) 是相邻阶之间的距离，不是峰的自身半高宽度，二者的比才刻画精细度。
(3) 对 \(R\to1\) 和 \(\delta=2\pi m+\epsilon\)，有 \(T\simeq[1+F\epsilon^2/4]^{-1}\)。半高条件为 \(|\epsilon|=2/\sqrt F\)，相位全宽 \(4/\sqrt F\)；利用 \(d\delta/d\nu=4\pi nL/c\) 得 \(\delta\nu_{\rm FWHM}\simeq\Delta\nu_{\rm FSR}(1-R)/(\pi\sqrt R)\)。精细度 \(\mathcal F=\Delta\nu_{\rm FSR}/\delta\nu_{\rm FWHM}\simeq\pi\sqrt R/(1-R)\)，谐振阶的 \(Q_{\rm cav}=\nu_m/\delta\nu_{\rm FWHM}\simeq m\mathcal F\)。
(4) 镜吸收使有效往返振幅乘积减小且峰值透射不再为一；腔内色散使频率间隔由群折射率 \(n_g=n+\nu\,dn/d\nu\) 控制，局部 \(\Delta\nu_{\rm FSR}\simeq c/(2n_gL)\)。提高 \(R\) 可缩窄单峰线宽但不会扩大邻峰的自由光谱间隔，因此强精细度仍可能把未知谱线混淆到相邻阶，需要额外的阶次标定。
